\documentclass[12pt,openany,a4paper]{article}
\usepackage{graphicx} 
\usepackage{multirow}
\usepackage{enumitem}
\usepackage{amssymb}
\usepackage{amsmath}
\usepackage{amsthm}
\usepackage{xcolor}
\usepackage{color, colortbl}
\usepackage{hyperref}
\hypersetup{
    colorlinks=true,
    linkcolor=blue,
    filecolor=magenta,      
    urlcolor=blue,
}
\usepackage{soul}
\usepackage{cancel}
\usepackage{geometry}
	\geometry{
        left=10mm,
        right=22mm,
        top=20mm,
        bottom=20mm,
    }
\usepackage{fancyhdr}
\usepackage{tikz, pgfplots, tkz-euclide,calc}
    \usetikzlibrary{patterns,snakes,shapes.arrows}

\newtheorem*{teorema}{Teorema}

\newcommand{\R}{\mathbb{R}}
\newcommand{\N}{\mathbb{N}}
\newcommand{\jawab}{\textbf{Jawab}:}
\renewcommand{\headrulewidth}{0pt}
\pagestyle{fancy}

\begin{document}
\fancyfoot[C]{\raisebox{.5ex}{\rule{0.5cm}{.4pt}}o0o\raisebox{.5ex}{\rule{0.5cm}{.4pt}}}

\noindent\textbf{\Large Kuis 2 Metode Matematika (A)}
\vspace{0.5cm}

\begin{enumerate}
  \item Dapatkan invers Transformasi $Z$ (barisan $y_n$) jika diberikan
        \[F(z) = \frac{z}{(z+1)(z+3)}.\]
        \textit{Petunjuk}: Gunakan tabel Transformasi $Z$ berikut

        \begin{table}[h]
          \centering
          \renewcommand{\arraystretch}{1.5}
          \caption{Beberapa Pasangan Transformasi $Z$}
          \begin{tabular}{clc}
            \hline
            No. & Barisan $f(n)$ & Transformasi $Z$, $F(z)$  \\
            \hline
            1   & $1$            & $\dfrac{z}{z-1}$          \\[1ex]
            2   & $n$            & $\dfrac{z}{(z-1)^2}$      \\[1ex]
            3   & $n^2$          & $\dfrac{z(z+1)}{(z-1)^3}$ \\[1ex]
            4   & $a^n$          & $\dfrac{z}{z-a}$          \\[1ex]
            5   & $n a^n$        & $\dfrac{az}{(z-a)^2}$     \\[1ex]
            \hline
          \end{tabular}
        \end{table}

  \item Selesaikan persamaan diferensial berikut dengan metode deret pangkat di sekitar $x = 1$. Tentukan solusi dalam bentuk deret hingga suku $(x - 1)^5$.
        \[y'' - (x - 1)y' = 0.\]

  \item Kembangkan fungsi berikut ke dalam deret polinomial Legendre pada interval $[-1, 1]$. Tentukan tiga suku pertama yang tidak nol.
        \[f(x) = \begin{cases} 2, & -1 \le x < 0, \\ 1, & 0 \le x < 1. \end{cases}\]
        \textit{Petunjuk}: Gunakan rumus koefisien deret Legendre
        \[a_n = \frac{2n + 1}{2} \int_{-1}^{1} f(x)P_n(x) \, dx,\]
        dengan beberapa polinomial Legendre pertama
        \[P_0(x) = 1, \quad P_1(x) = x, \quad P_2(x) = \frac{1}{2}(3x^2 - 1), \quad P_3(x) = \frac{1}{2}(5x^3 - 3x).\]

  \item Dengan menggunakan integrasi parsial berulang dan rumus rekurensi fungsi Bessel
        \[\frac{d}{dx} (x^n J_n(x)) = x^n J_{n-1}(x),\]
        buktikan bahwa
        \[\int x^5 J_2(x) \, dx = x^5 J_3(x) - 2x^4 J_4(x) + C.\]
        \textit{Petunjuk}: Gunakan integrasi parsial
        \[\int u \, dv = uv - \int v \, du,\]
        dengan
        \[u = x^2, \quad dv = x^3 J_2(x) \, dx.\]
\end{enumerate}

\newpage
\textbf{SOLUSI}
\begin{enumerate}
  \item Kita akan mencari invers transformasi $Z$ dari fungsi
        \[F(z) = \frac{z}{(z+1)(z+3)}.\]
        Kita dapat menuliskan kembali fungsi tersebut menjadi
        \[\frac{F(z)}{z} = \frac{1}{(z+1)(z+3)}.\]
        Dengan menggunakan dekomposisi pecahan parsial,
        \[\frac{1}{(z+1)(z+3)} = \frac{A}{z+1} + \frac{B}{z+3}.\]
        Maka $A(z+3) + B(z+1) = 1$.
        \begin{itemize}
          \item Substitusi $z = -1 \implies A(-1+3) = 1 \implies 2A = 1 \implies A = \frac{1}{2}$.
          \item Substitusi $z = -3 \implies B(-3+1) = 1 \implies -2B = 1 \implies B = -\frac{1}{2}$.
        \end{itemize}
        Sehingga kita peroleh
        \[\frac{F(z)}{z} = \frac{\frac{1}{2}}{z+1} - \frac{\frac{1}{2}}{z+3}\]
        \[F(z) = \frac{1}{2}\frac{z}{z-(-1)} - \frac{1}{2}\frac{z}{z-(-3)}.\]
        Berdasarkan tabel transformasi $Z$ nomor 4, kita mengetahui bahwa $a^n \leftrightarrow \dfrac{z}{z-a}$.
        Dengan demikian, invers transformasi $Z$ dari $\dfrac{z}{z-(-1)}$ adalah $(-1)^n$ dan invers dari $\dfrac{z}{z-(-3)}$ adalah $(-3)^n$.
        Jadi, barisan $y_n$ adalah
        \[y_n = \frac{1}{2}(-1)^n - \frac{1}{2}(-3)^n, \quad \text{untuk } n \ge 0.\]

  \item Diberikan persamaan diferensial
        \[y'' - (x - 1)y' = 0.\]
        Kita akan menyelesaikannya dengan metode deret pangkat di sekitar $x = 1$.
        Misalkan $t = x - 1$, maka $y(x) = y(t+1) = Y(t)$. Persamaan menjadi
        \[Y''(t) - tY'(t) = 0.\]
        Asumsikan solusi berupa deret pangkat:
        \[Y(t) = \sum_{n=0}^{\infty} c_n t^n.\]
        Maka turunan pertama dan keduanya adalah:
        \[Y'(t) = \sum_{n=1}^{\infty} n c_n t^{n-1},\]
        \[Y''(t) = \sum_{n=2}^{\infty} n(n-1) c_n t^{n-2}.\]
        Substitusi deret-deret ini ke dalam persamaan diferensial $Y'' - tY' = 0$:
        \[\sum_{n=2}^{\infty} n(n-1) c_n t^{n-2} - t\sum_{n=1}^{\infty} n c_n t^{n-1} = 0,\]
        \[\sum_{n=2}^{\infty} n(n-1) c_n t^{n-2} - \sum_{n=1}^{\infty} n c_n t^n = 0.\]
        Agar pangkat dari $t$ sama, kita geser indeks pada jumlahan pertama dengan memisalkan $k = n-2 \implies n = k+2$:
        \[\sum_{k=0}^{\infty} (k+2)(k+1) c_{k+2} t^k - \sum_{k=1}^{\infty} k c_k t^k = 0.\]
        Kita pisahkan suku untuk $k=0$:
        \[(2)(1) c_2 + \sum_{k=1}^{\infty} \left[ (k+2)(k+1) c_{k+2} - k c_k \right] t^k = 0.\]
        Dari persamaan di atas, dengan menyamakan koefisien untuk setiap pangkat $t$, kita peroleh:
        \begin{itemize}
          \item Untuk konstanta ($k=0$): $2 c_2 = 0 \implies c_2 = 0$.
          \item Untuk $k \ge 1$: $(k+2)(k+1) c_{k+2} - k c_k = 0 \implies c_{k+2} = \frac{k}{(k+2)(k+1)} c_k$.
        \end{itemize}
        Kita akan mencari nilai koefisien $c_n$ hingga pangkat $t^5$:
        \begin{itemize}
          \item $k = 1$: $c_3 = \frac{1}{(3)(2)} c_1 = \frac{1}{6} c_1$.
          \item $k = 2$: $c_4 = \frac{2}{(4)(3)} c_2 = 0$ (karena $c_2 = 0$).
          \item $k = 3$: $c_5 = \frac{3}{(5)(4)} c_3 = \frac{3}{20} \left(\frac{1}{6} c_1\right) = \frac{1}{40} c_1$.
        \end{itemize}
        Solusi deret hingga suku $t^5$ adalah:
        \[Y(t) = c_0 + c_1 t + c_2 t^2 + c_3 t^3 + c_4 t^4 + c_5 t^5 + \dots\]
        \[Y(t) = c_0 + c_1 t + \frac{1}{6} c_1 t^3 + \frac{1}{40} c_1 t^5 + \dots\]
        Dengan mengembalikan $t = x - 1$, kita dapatkan solusi $y(x)$:
        \[y(x) = c_0 + c_1 (x-1) + \frac{1}{6} c_1 (x-1)^3 + \frac{1}{40} c_1 (x-1)^5 + \dots\]
        dengan $c_0, c_1$ adalah konstanta sebarang.

  \item Fungsi $f(x)$ diberikan oleh
        \[f(x) = \begin{cases} 2, & -1 \le x < 0, \\ 1, & 0 \le x < 1. \end{cases}\]
        Kita akan mengembangkan fungsi ini ke dalam deret polinomial Legendre:
        \[f(x) = \sum_{n=0}^{\infty} a_n P_n(x).\]
        Koefisien deret dihitung dengan rumus:
        \[a_n = \frac{2n + 1}{2} \int_{-1}^{1} f(x)P_n(x) \, dx.\]
        Kita akan menghitung koefisien-koefisien awal untuk mendapatkan tiga suku pertama yang tidak nol.
        \begin{itemize}
          \item \textbf{Suku ke-0 ($n=0$)}:\\
                $P_0(x) = 1$.
                \begin{align*}
                  a_0 & = \frac{1}{2} \int_{-1}^{1} f(x)(1) \, dx = \frac{1}{2} \left[ \int_{-1}^{0} 2 \, dx + \int_{0}^{1} 1 \, dx \right] \\
                      & = \frac{1}{2} \left[ 2(x)\Big|_{-1}^0 + x\Big|_0^1 \right] = \frac{1}{2} \left[ 2(0 - (-1)) + (1 - 0) \right]       \\
                      & = \frac{1}{2} [2 + 1] = \frac{3}{2}.
                \end{align*}
                Suku pertama yang tidak nol: $\frac{3}{2}P_0(x)$.

          \item \textbf{Suku ke-1 ($n=1$)}:\\
                $P_1(x) = x$.
                \begin{align*}
                  a_1 & = \frac{3}{2} \int_{-1}^{1} f(x)(x) \, dx = \frac{3}{2} \left[ \int_{-1}^{0} 2x \, dx + \int_{0}^{1} x \, dx \right]                       \\
                      & = \frac{3}{2} \left[ x^2\Big|_{-1}^0 + \frac{1}{2}x^2\Big|_0^1 \right] = \frac{3}{2} \left[ (0 - 1) + \left(\frac{1}{2} - 0\right) \right] \\
                      & = \frac{3}{2} \left[ -1 + \frac{1}{2} \right] = \frac{3}{2} \left[ -\frac{1}{2} \right] = -\frac{3}{4}.
                \end{align*}
                Suku kedua yang tidak nol: $-\frac{3}{4}P_1(x)$.

          \item \textbf{Suku ke-2 ($n=2$)}:\\
                $P_2(x) = \frac{1}{2}(3x^2 - 1)$.
                \begin{align*}
                  a_2 & = \frac{5}{2} \int_{-1}^{1} f(x) P_2(x) \, dx = \frac{5}{2} \left[ \int_{-1}^{0} 2 P_2(x) \, dx + \int_{0}^{1} P_2(x) \, dx \right]
                \end{align*}
                Kita ketahui bahwa $P_2(x)$ adalah fungsi genap, sehingga $\int_{-1}^{0} P_2(x) \, dx = \int_{0}^{1} P_2(x) \, dx$.
                Mari kita hitung integral $\int_{0}^{1} P_2(x) \, dx$:
                \[\int_{0}^{1} \frac{1}{2}(3x^2 - 1) \, dx = \frac{1}{2} \left[ x^3 - x \right]_0^1 = \frac{1}{2} (1 - 1) = 0.\]
                Karena integral di selang $[0,1]$ bernilai $0$, maka integral di selang $[-1,0]$ juga $0$. Sehingga:
                \[a_2 = \frac{5}{2} [ 2(0) + 0 ] = 0.\]

          \item \textbf{Suku ke-3 ($n=3$)}:\\
                $P_3(x) = \frac{1}{2}(5x^3 - 3x)$.
                \begin{align*}
                  a_3 & = \frac{7}{2} \int_{-1}^{1} f(x) P_3(x) \, dx = \frac{7}{2} \left[ \int_{-1}^{0} 2 P_3(x) \, dx + \int_{0}^{1} P_3(x) \, dx \right]
                \end{align*}
                Fungsi $P_3(x)$ adalah fungsi ganjil, sehingga $\int_{-1}^{0} P_3(x) \, dx = -\int_{0}^{1} P_3(x) \, dx$.
                Mari kita hitung integral $\int_{0}^{1} P_3(x) \, dx$:
                \[\int_{0}^{1} \frac{1}{2}(5x^3 - 3x) \, dx = \frac{1}{2} \left[ \frac{5}{4}x^4 - \frac{3}{2}x^2 \right]_0^1 = \frac{1}{2} \left( \frac{5}{4} - \frac{6}{4} \right) = \frac{1}{2} \left( -\frac{1}{4} \right) = -\frac{1}{8}.\]
                Maka,
                \[\int_{-1}^{0} P_3(x) \, dx = -\left(-\frac{1}{8}\right) = \frac{1}{8}.\]
                Sehingga:
                \begin{align*}
                  a_3 & = \frac{7}{2} \left[ 2\left(\frac{1}{8}\right) + \left(-\frac{1}{8}\right) \right] = \frac{7}{2} \left[ \frac{2}{8} - \frac{1}{8} \right] = \frac{7}{2} \left[ \frac{1}{8} \right] = \frac{7}{16}.
                \end{align*}
                Suku ketiga yang tidak nol: $\frac{7}{16}P_3(x)$.
        \end{itemize}
        Jadi, ekspansi deret polinomial Legendre hingga tiga suku pertama yang tidak nol adalah:
        \[f(x) \approx \frac{3}{2}P_0(x) - \frac{3}{4}P_1(x) + \frac{7}{16}P_3(x).\]

  \item Akan dibuktikan bahwa
        \[\int x^5 J_2(x) \, dx = x^5 J_3(x) - 2x^4 J_4(x) + C.\]
        Kita gunakan rumus rekurensi $\frac{d}{dx} (x^n J_n(x)) = x^n J_{n-1}(x)$. Dari sini, kita peroleh
        \[\int x^n J_{n-1}(x) \, dx = x^n J_n(x).\]
        Gunakan metode integrasi parsial $\int u \, dv = uv - \int v \, du$ dengan pemisalan:
        \[u = x^2 \quad \implies \quad du = 2x \, dx\]
        \[dv = x^3 J_2(x) \, dx \quad \implies \quad v = \int x^3 J_2(x) \, dx.\]
        Berdasarkan rumus integral rekurensi dengan $n=3$, kita dapatkan:
        \[v = \int x^3 J_{3-1}(x) \, dx = x^3 J_3(x).\]
        Maka integrasi parsialnya menjadi:
        \begin{align*}
          \int x^5 J_2(x) \, dx & = \int (x^2) \left( x^3 J_2(x) \, dx \right)                 \\
                                & = \int u \, dv                                               \\
                                & = uv - \int v \, du                                          \\
                                & = x^2 \left( x^3 J_3(x) \right) - \int x^3 J_3(x) (2x \, dx) \\
                                & = x^5 J_3(x) - 2 \int x^4 J_3(x) \, dx.
        \end{align*}
        Selanjutnya, untuk menyelesaikan integral $\int x^4 J_3(x) \, dx$, kita gunakan kembali rumus integral rekurensi untuk $n=4$:
        \[\int x^4 J_{4-1}(x) \, dx = x^4 J_4(x).\]
        Maka persamaan sebelumnya menjadi:
        \begin{align*}
          \int x^5 J_2(x) \, dx & = x^5 J_3(x) - 2 \left( x^4 J_4(x) \right) + C \\
                                & = x^5 J_3(x) - 2x^4 J_4(x) + C.
        \end{align*}
        Terbukti.
\end{enumerate}

\end{document}
